\chapter{Database Relationships}

\section{One-to-One: User and Profile}

\begin{diagrambox}[One-to-One]
\begin{verbatim}
User { _id: "u1", name: "Asha" }
Profile { _id: "p1", user: "u1" (ref User), bio: "...", location: "Kolkata" }
\end{verbatim}
\end{diagrambox}

\begin{lstlisting}[language=JavaScript]
// models/Profile.js
const profileSchema = new mongoose.Schema({
  user: { type: mongoose.Schema.Types.ObjectId, ref: "User", required: true, unique: true },
  bio: String,
  location: String,
});
\end{lstlisting}

\section{One-to-Many: User and Tasks}

\begin{diagrambox}[One-to-Many]
\begin{verbatim}
User "u1"
  |-- Task "t1" { user: "u1", title: "Fix login bug" }
  |-- Task "t2" { user: "u1", title: "Write tests" }
  `-- Task "t3" { user: "u1", title: "Deploy to prod" }
\end{verbatim}
\end{diagrambox}

This is the \texttt{Task.user} field from Chapter 13: each Task stores an \texttt{ObjectId} pointing back to its owning User.

\section{populate(): Resolving References}

By default Mongoose returns only the raw \texttt{ObjectId} for a \texttt{ref} field. \texttt{.populate()} replaces it with the actual document.

\begin{lstlisting}[language=JavaScript]
// Without populate
const task = await Task.findById(taskId);

// With populate - only pull name and email
const task = await Task.find().populate("user", "name email");
\end{lstlisting}

\subsection{Before / After populate("user", "name email")}

\begin{lstlisting}[language=JavaScript]
// Before: user: "664e0a1b2c3d4e5f6a7b8c9d"
// After:
{
  "title": "Fix login bug",
  "status": "todo",
  "user": { "_id": "664e0a1b2c3d4e5f6a7b8c9d", "name": "Asha Verma", "email": "asha@mail.com" }
}
\end{lstlisting}

\begin{notebox}[NOTE]
The \texttt{"name email"} argument is a projection limiting which fields come back, keeping payloads small and avoiding exposing fields like \texttt{password}.
\end{notebox}

\section{Many-to-Many: Users and Projects}

A many-to-many relation needs an array of references on at least one side, here a Project's array of member User ids.

\begin{lstlisting}[language=JavaScript]
// models/Project.js
const projectSchema = new mongoose.Schema(
  {
    name: { type: String, required: true },
    members: [{ type: mongoose.Schema.Types.ObjectId, ref: "User" }],
  },
  { timestamps: true }
);
\end{lstlisting}

\begin{diagrambox}[Many-to-Many]
\begin{verbatim}
Project "Website Redesign" members: [u1, u2, u3]
Project "Mobile App"       members: [u2, u4]
u2 belongs to both projects.
\end{verbatim}
\end{diagrambox}

\begin{lstlisting}[language=JavaScript]
const project = await Project.findById(projectId).populate("members", "name email");
// project.members is now an array of user objects, not ids
\end{lstlisting}

\begin{tipbox}[TIP]
\texttt{.populate()} works the same for a single \texttt{ObjectId} or an array --- Mongoose detects the schema type automatically.
\end{tipbox}

\begin{warnbox}[WARNING]
Populate runs an extra query; only populate fields the response actually needs to avoid hurting performance.
\end{warnbox}
